\documentclass[11pt]{article}
\usepackage[a4paper,margin=2.1cm]{geometry}
\usepackage[utf8]{inputenc}
\usepackage[spanish]{babel}
\usepackage[T1]{fontenc}
\usepackage{lmodern}
\usepackage{microtype}
\sloppy
\emergencystretch=3em
\setlength{\headheight}{14.5pt}
\usepackage{xcolor}
\definecolor{vino}{RGB}{150,25,40}
\definecolor{grisclaro}{RGB}{242,242,242}
\usepackage{titlesec}
\titleformat{\section}{\normalfont\Large\bfseries\color{vino}}{\thesection.}{0.5em}{}
\titleformat{\subsection}{\normalfont\large\bfseries\color{vino}}{\thesubsection}{0.5em}{}
\usepackage{listings}
\usepackage{booktabs}
\usepackage{array}
\usepackage{longtable}
\usepackage{enumitem}
\usepackage{hyperref}
\hypersetup{colorlinks=true, linkcolor=vino, urlcolor=vino, citecolor=vino}
\usepackage{fancyhdr}
\pagestyle{fancy}
\fancyhf{}
\renewcommand{\headrulewidth}{0.4pt}
\fancyhead[L]{Laboratorio NoSQL --- MongoDB (M1721)}
\fancyhead[R]{Grupo 04}
\fancyfoot[C]{\thepage}
\setlength{\parskip}{4pt}
\setlength{\parindent}{0pt}

\lstdefinelanguage{JSMongo}{
  morekeywords={const,function,print,printjson,forEach,db},
  sensitive=true,
  morecomment=[l]{//},
  morestring=[b]",
}
\lstset{
  language=JSMongo,
  basicstyle=\ttfamily\scriptsize,
  keywordstyle=\color{vino}\bfseries,
  commentstyle=\color{gray}\itshape,
  stringstyle=\color{teal},
  breaklines=true,
  breakatwhitespace=false,
  columns=fullflexible,
  backgroundcolor=\color{grisclaro},
  frame=single,
  rulecolor=\color{gray!40},
  xleftmargin=2pt, xrightmargin=2pt,
  showstringspaces=false,
  tabsize=2,
}

\title{\vspace{-1.5cm}\textbf{\Huge Laboratorio NoSQL --- MongoDB}\\[0.3cm]
\Large Diseño documental y consultas sobre evidencia VIIRS de fuego (celda-mes)\\[0.2cm]
\normalsize Asignatura M1721 --- Bases de datos para Datos Estructurados, no Estructurados\\
e Introducción a Big Data para Ciencia de Datos}
\author{Grupo 04 \\ Proyecto FireForest (componente NoSQL independiente)}
\date{Quito/Loja, 2026}

\begin{document}
\maketitle
\thispagestyle{empty}
\vspace{1cm}
\begin{center}
\begin{tabular}{lp{9cm}}
\toprule
\textbf{Base de datos independiente} & \texttt{taller\_nosql\_fireforest} \\
\textbf{Colección principal} & \texttt{evidencia\_viirs\_celda\_mes} \\
\textbf{Documentos cargados} & 105 (datos reales de VIIRS, no simulados) \\
\textbf{Fuente} & Registros reales de VIIRS procesados previamente (solo lectura), v1.0.3, 2019--2025 \\
\bottomrule
\end{tabular}
\end{center}
\newpage

\section{Descripción del caso y justificación de MongoDB}

Los incendios forestales del cantón Loja generan información heterogénea:
detecciones satelitales VIIRS agregadas por celda espacial de 500\,m y mes,
indicadores de calidad de observación, evidencia de fuego bajo distintos
criterios de confianza, y metadatos operativos que se agregan \emph{después}
de la carga inicial (revisiones, bitácora de procesamiento). Esta
heterogeneidad —estructuras anidadas, un arreglo de evidencia de cardinalidad
fija, y campos que aparecen solo después de una actualización— es más
natural de representar como documentos que como tablas relacionales rígidas:
un único \lstinline!findOne()! recupera toda la unidad celda-mes sin
necesidad de \texttt{JOIN}, y los metadatos operativos se incorporan con
\lstinline!\$set!/\lstinline!\$push! sin alterar el esquema.

Este laboratorio reutiliza la unidad de análisis celda-mes del Proyecto
Integrador FireForest (ver \texttt{AGENTS.md}) y constituye su componente
NoSQL. Preguntas que responde: (1) ¿qué registros celda-mes superan un
umbral de FRP relevante y cuándo ocurrieron? (2) ¿en qué meses y años se
concentra la mayor intensidad térmica, y cómo difiere la evidencia de fuego
entre el criterio amplio y el criterio nominal/alta de confianza?

\section{Colección principal y unidad documental}

\textbf{Colección principal:} \texttt{evidencia\_viirs\_celda\_mes}, dentro
de la base independiente \texttt{taller\_nosql\_fireforest} (distinta de la
base \texttt{fireforest} del Proyecto Integrador, que no fue modificada).

\textbf{Unidad documental:} cada documento representa \textbf{una celda
espacial de 500\,m observada durante un mes calendario} (celda-mes),
identificada de forma única por \lstinline!celda.cell_index + periodo.anio + periodo.mes!
(reforzado por el índice único \texttt{ux\_celda\_anio\_mes}).

\section{Esquema del documento}

\begin{lstlisting}[language={}, basicstyle=\ttfamily\footnotesize]
evidencia_viirs_celda_mes
|- _id                          string
|- celda            { cell_index, unidad_espacial }
|- periodo          { anio, mes, anio_mes, fecha_inicio: Date }
|- fuente_satelital { fuente_primaria, asset_gee, procedencia, version_origen }
|- metricas
|   |- detecciones  { media_todas, media_nominal_alta }
|   |- presencia    { fraccion_todas, fraccion_nominal_alta }
|   \- frp          { suma_observada_media_mw, maxima_media_mw }
|- calidad_observacion { dias_validos_media, disponibilidad_pct, ... }
|- evidencias [ { categoria, detecciones_media, presencia_fraccion,
|                 evidencia_fuego: bool }, ... ]   (2 elementos)
|- muestreo         { estrato, metodo, semilla }
|- metadatos        { unidad_analisis, datos_simulados, uso, estado_revision? }
|- revisado_manualmente?        bool
\- bitacora_laboratorio? [ { accion, responsable, fecha_procesamiento: Date } ]
\end{lstlisting}

\section{Decisiones de embebido y referencia}

Se \textbf{embebieron} \lstinline!celda!, \lstinline!periodo!,
\lstinline!fuente_satelital!, \lstinline!metricas! y
\lstinline!calidad_observacion! porque tienen cardinalidad 1--1, se leen
siempre junto con el documento y no cambian tras la carga. El arreglo
\lstinline!evidencias! se embebió por tener cardinalidad fija y pequeña (2
elementos: \texttt{todas} y \texttt{nominal\_alta}), consultado con
\lstinline!$elemMatch! (consultas 2 y 6) o aplanado puntualmente con
\lstinline!$unwind! (consulta 9). Los metadatos operativos
(\lstinline!metadatos.estado_revision!, \lstinline!revisado_manualmente!,
\lstinline!bitacora_laboratorio!) también se embebieron pese a evolucionar
después de la carga, porque su cardinalidad sigue siendo pequeña y ligada al
documento.

\textbf{No se usaron referencias ni \lstinline!$lookup!}: no existe en este
caso una entidad de cardinalidad alta e independiente que se consulte por
separado del documento celda-mes; las 10 consultas exigidas se resuelven
íntegramente sobre esta única colección. Una bitácora de revisión más
extensa (con múltiples analistas y comentarios largos) sí ameritaría una
colección referenciada aparte; no es el caso de este taller.

\section{Documento JSON/BSON completo (ejemplo real)}

\begin{lstlisting}
{
  "_id": "VCM_0001",
  "celda": { "cell_index": 1803, "unidad_espacial": "celda_500m" },
  "periodo": { "anio": 2019, "mes": 1, "anio_mes": "2019-01",
               "fecha_inicio": ISODate("2019-01-01T00:00:00Z") },
  "fuente_satelital": { "fuente_primaria": "NASA FIRMS - VIIRS",
    "asset_gee": "VIIRS_PREP_2019_01",
    "procedencia": "registros_previamente_procesados",
    "version_origen": "v1.0.3" },
  "metricas": {
    "detecciones": { "media_todas": 1, "media_nominal_alta": 0 },
    "presencia": { "fraccion_todas": 1, "fraccion_nominal_alta": 0 },
    "frp": { "suma_observada_media_mw": 4.7, "maxima_media_mw": 4.7 }
  },
  "calidad_observacion": { "dias_validos_media": 31,
    "disponibilidad_pct": 100, "observado_media": 1 },
  "evidencias": [
    { "categoria": "todas", "detecciones_media": 1,
      "presencia_fraccion": 1, "evidencia_fuego": true },
    { "categoria": "nominal_alta", "detecciones_media": 0,
      "presencia_fraccion": 0, "evidencia_fuego": false }
  ],
  "muestreo": { "estrato": "con_fuego",
    "metodo": "estratificado_por_anio_y_evidencia", "semilla": 2026 },
  "metadatos": { "unidad_analisis": "celda_mes_500m",
    "datos_simulados": false, "uso": "Laboratorio NoSQL MongoDB - M1721" }
}
\end{lstlisting}

\section{Evidencia de carga de los 105 documentos}

Verificado mediante ejecución real de
\texttt{07\_entrega/Grupo04\_Laboratorio\_NoSQL\_MongoDB.js} contra
\texttt{taller\_nosql\_fireforest} (detalle completo en
\texttt{05\_resultados/01\_validacion\_carga.txt}):

\begin{center}
\begin{tabular}{lc}
\toprule
\textbf{Verificación} & \textbf{Resultado real} \\
\midrule
Conteo total de documentos & 105 \\
\texttt{\_id} únicos & 105 \\
Combinaciones celda-año-mes únicas & 105 \\
Fechas BSON (\texttt{periodo.fecha\_inicio}) & 105 \\
Objetos anidados (\texttt{metricas}) & 105 \\
Arreglos (\texttt{evidencias}) & 105 \\
Booleanos (\texttt{metadatos.datos\_simulados}) & 105 \\
Con evidencia de fuego (criterio \texttt{todas}) & 70 \\
Sin evidencia de fuego & 35 \\
Con evidencia nominal/alta & 68 \\
FRP máxima de la muestra & 122,3162 MW (VCM\_0073, 2023-10) \\
\bottomrule
\end{tabular}
\end{center}

\textbf{Índices creados:} \texttt{ux\_celda\_anio\_mes} (único, sobre
\lstinline!celda.cell_index+periodo.anio+periodo.mes!), \texttt{ix\_fecha\_inicio},
\texttt{ix\_frp\_maxima} (descendente) y \texttt{ix\_evidencias\_categoria\_fuego}
(multiclave, sobre \lstinline!evidencias.categoria! y
\lstinline!evidencias.evidencia_fuego!).

\section{Consultas MongoDB (Parte C) --- resumen de las 10 consultas}

Todas las consultas incluyen orden secundario por \texttt{\_id} (o la clave
de agrupación) para evitar ambigüedad ante empates. Código completo,
resultado sin recortar e interpretación de cada consulta:
\texttt{04\_consultas/01\_consultas\_1\_a\_10.md} y
\texttt{05\_resultados/02\_resultado\_consultas\_1\_10.txt}.

\begin{center}
\small
\begin{longtable}{p{0.4cm}p{2.6cm}p{4.6cm}p{5.6cm}}
\toprule
\textbf{\#} & \textbf{Tipo} & \textbf{Pregunta} & \textbf{Resultado real (clave)} \\
\midrule
1 & Filtro \lstinline!$gt! & FRP máxima $>$ 20 MW? & 13 registros; máximo VCM\_0073 (122{,}32 MW) \\\midrule
2 & Filtro combinado & 2023-2024 con FRP $\geq$10 y evid. nominal/alta? & 11 registros cumplen ambos criterios \\\midrule
3 & Proyección & Campos esenciales de 10 registros con fuego & 10 documentos, solo 5 campos c/u \\\midrule
4 & Orden y límite & 5 registros con mayor FRP máxima & VCM\_0073 (122{,}32), VCM\_0009 (90{,}0), \ldots \\\midrule
5 & Campos anidados & Disponibilidad 100\% y $\geq$1 detección media & 10 registros cumplen ambas condiciones \\\midrule
6 & Arreglos \lstinline!$elemMatch! & Evidencia nominal/alta con detección $>$0 & 10 registros \\\midrule
7 & Pipeline mensual & Mes con mayor intensidad térmica & Sep.: 350{,}83 MW acum.; Oct.: 30{,}58 MW prom., máx. 122{,}32 MW \\\midrule
8 & Pipeline anual & Variación anual de métricas de fuego & 2023: mayor FRP acum. (349{,}58) y prom. (34{,}96 MW) \\\midrule
9 & \lstinline!$unwind! & Evidencia total vs. nominal/alta & 70 vs. 68 documentos con fuego, sobre 105 evaluados \\\midrule
10 & Integradora & Dataset plano celda-mes para análisis/ML & 15 filas; encabeza VCM\_0073 (188{,}27 MW) \\\midrule
\caption{Resumen de las 10 consultas exigidas por la guía, con resultado real obtenido.}
\end{longtable}
\end{center}

\subsection{Código completo de las 10 consultas}

\begin{lstlisting}
// Consulta 1 (filtro $gt)
coleccion.find(
  { "metricas.frp.maxima_media_mw": { $gt: 20 } },
  { _id: 1, "celda.cell_index": 1, "periodo.anio_mes": 1,
    "metricas.frp.maxima_media_mw": 1 }
).sort({ "metricas.frp.maxima_media_mw": -1, _id: 1 }).forEach(printjson);

// Consulta 2 (filtro combinado $in + $gte + $elemMatch)
coleccion.find(
  { "periodo.anio": { $in: [2023, 2024] },
    "metricas.frp.maxima_media_mw": { $gte: 10 },
    evidencias: { $elemMatch: { categoria: "nominal_alta", evidencia_fuego: true } } },
  { _id: 1, "periodo.anio_mes": 1, "celda.cell_index": 1,
    "metricas.frp.maxima_media_mw": 1 }
).sort({ "periodo.fecha_inicio": 1, _id: 1 }).forEach(printjson);

// Consulta 3 (proyeccion)
coleccion.find(
  { "metricas.frp.suma_observada_media_mw": { $gt: 0 } },
  { _id: 1, "celda.cell_index": 1, "periodo.anio_mes": 1,
    "metricas.detecciones.media_todas": 1,
    "metricas.frp.suma_observada_media_mw": 1 }
).sort({ "periodo.fecha_inicio": 1, _id: 1 }).limit(10).forEach(printjson);

// Consulta 4 (orden y limite)
coleccion.find(
  {}, { _id: 1, "celda.cell_index": 1, "periodo.anio_mes": 1,
        "metricas.frp.maxima_media_mw": 1 }
).sort({ "metricas.frp.maxima_media_mw": -1, _id: 1 }).limit(5).forEach(printjson);

// Consulta 5 (campos anidados)
coleccion.find(
  { "calidad_observacion.disponibilidad_pct": 100,
    "metricas.detecciones.media_todas": { $gte: 1 } },
  { _id: 1, "periodo.anio_mes": 1, "celda.cell_index": 1,
    "calidad_observacion.disponibilidad_pct": 1,
    "metricas.detecciones.media_todas": 1 }
).sort({ "metricas.detecciones.media_todas": -1, _id: 1 }).limit(10).forEach(printjson);

// Consulta 6 (arreglos con $elemMatch)
coleccion.find(
  { evidencias: { $elemMatch: { categoria: "nominal_alta",
      evidencia_fuego: true, detecciones_media: { $gt: 0 } } } },
  { _id: 1, "celda.cell_index": 1, "periodo.anio_mes": 1,
    evidencias: { $elemMatch: { categoria: "nominal_alta" } } }
).sort({ "periodo.fecha_inicio": 1, _id: 1 }).limit(10).forEach(printjson);

// Consulta 7 ($match + $group + $sort, por mes calendario)
coleccion.aggregate([
  { $match: { "metricas.frp.suma_observada_media_mw": { $gt: 0 } } },
  { $group: { _id: "$periodo.mes", registros_con_fuego: { $sum: 1 },
      frp_total_muestra_mw: { $sum: "$metricas.frp.suma_observada_media_mw" },
      frp_promedio_mw: { $avg: "$metricas.frp.suma_observada_media_mw" },
      frp_maxima_mw: { $max: "$metricas.frp.maxima_media_mw" } } },
  { $project: { _id: 0, mes: "$_id", registros_con_fuego: 1,
      frp_total_muestra_mw: { $round: ["$frp_total_muestra_mw", 2] },
      frp_promedio_mw: { $round: ["$frp_promedio_mw", 2] },
      frp_maxima_mw: { $round: ["$frp_maxima_mw", 2] } } },
  { $sort: { frp_total_muestra_mw: -1, mes: 1 } }
]).forEach(printjson);

// Consulta 8 ($match + $group + $sort, por anio)
coleccion.aggregate([
  { $match: { "metricas.frp.suma_observada_media_mw": { $gt: 0 } } },
  { $group: { _id: "$periodo.anio", registros_con_fuego: { $sum: 1 },
      frp_promedio_mw: { $avg: "$metricas.frp.suma_observada_media_mw" },
      frp_total_muestra_mw: { $sum: "$metricas.frp.suma_observada_media_mw" },
      frp_maxima_mw: { $max: "$metricas.frp.maxima_media_mw" } } },
  { $project: { _id: 0, anio: "$_id", registros_con_fuego: 1,
      frp_promedio_mw: { $round: ["$frp_promedio_mw", 2] },
      frp_total_muestra_mw: { $round: ["$frp_total_muestra_mw", 2] },
      frp_maxima_mw: { $round: ["$frp_maxima_mw", 2] } } },
  { $sort: { frp_total_muestra_mw: -1, anio: 1 } }
]).forEach(printjson);

// Consulta 9 ($unwind + $group, evidencia total vs. nominal/alta)
coleccion.aggregate([
  { $unwind: "$evidencias" },
  { $group: { _id: "$evidencias.categoria", documentos_evaluados: { $sum: 1 },
      documentos_con_fuego: { $sum: { $cond: ["$evidencias.evidencia_fuego", 1, 0] } },
      detecciones_promedio: { $avg: "$evidencias.detecciones_media" } } },
  { $project: { _id: 0, categoria: "$_id", documentos_evaluados: 1,
      documentos_con_fuego: 1,
      detecciones_promedio: { $round: ["$detecciones_promedio", 3] } } },
  { $sort: { categoria: 1 } }
]).forEach(printjson);

// Consulta 10 (integradora: dataset plano para analisis/ML)
coleccion.aggregate([
  { $match: { "metricas.frp.suma_observada_media_mw": { $gt: 0 },
      "calidad_observacion.disponibilidad_pct": { $gte: 95 } } },
  { $set: { frp_por_deteccion_mw: { $cond: [
        { $gt: ["$metricas.detecciones.media_todas", 0] },
        { $divide: ["$metricas.frp.suma_observada_media_mw",
                    "$metricas.detecciones.media_todas"] }, 0 ] } } },
  { $project: { _id: 0, registro_id: "$_id", anio_mes: "$periodo.anio_mes",
      frp_total_media_mw: { $round: ["$metricas.frp.suma_observada_media_mw", 2] },
      frp_por_deteccion_mw: { $round: ["$frp_por_deteccion_mw", 2] } } },
  { $sort: { frp_total_media_mw: -1, registro_id: 1 } },
  { $limit: 15 }
]).forEach(printjson);
\end{lstlisting}

\subsection{Interpretación de los dos pipelines analíticos principales}

\textbf{Consulta 7 (por mes calendario):} septiembre concentra la mayor FRP
acumulada de la muestra (350,83~MW en 22 registros), mientras que octubre
presenta la mayor FRP promedio (30,58~MW) y el máximo absoluto
(122,32~MW): más eventos en septiembre, pero más intensos en octubre.

\textbf{Consulta 8 (por año):} 2023 combina la mayor FRP acumulada
(349,58~MW) y la mayor FRP promedio (34,96~MW) entre los siete años de la
muestra; 2022 presenta los valores más bajos (58,74~MW acumulados, 5,87~MW
promedio).

\section{Operaciones de actualización (Parte D)}

Las tres actualizaciones agregan o modifican \textbf{únicamente metadatos
operativos del laboratorio}; ningún valor satelital real se modifica.
Evidencia completa antes/después en
\texttt{05\_resultados/03\_actualizaciones\_antes\_despues.txt}.

\begin{center}
\small
\begin{tabular}{p{2.4cm}p{2.4cm}p{5.3cm}p{2.6cm}}
\toprule
\textbf{Operación} & \textbf{Alcance} & \textbf{Efecto real} & \textbf{Resultado} \\
\midrule
\texttt{updateOne()},\newline campo simple &
1 documento &
Agrega \lstinline!revisado_manualmente: true! al documento de mayor FRP (VCM\_0073) &
matched:1, modified:1 \\
\midrule
\texttt{updateMany()},\newline campo anidado &
70 documentos &
Agrega \lstinline!metadatos.estado_revision: "revisado_taller_nosql"! a los documentos con fuego &
matched:70, modified:70 \\
\midrule
\texttt{updateMany()} + \lstinline!$push! &
105 documentos &
Incorpora el arreglo \lstinline!bitacora_laboratorio! (acción, responsable, fecha) &
matched:105, modified:105 \\
\bottomrule
\end{tabular}
\end{center}

\section{Unidad de análisis y conexión con Ciencia de Datos}

\textbf{Unidad de análisis:} una celda-mes
(\lstinline!celda.cell_index+periodo.anio+periodo.mes!), consistente con la
unidad declarada para todo el Proyecto Integrador FireForest.

\textbf{Variable objetivo:} \lstinline!metricas.frp.maxima_media_mw!
(continua, intensidad térmica máxima) o
\lstinline!evidencias[categoria="todas"].evidencia_fuego! (binaria,
presencia de fuego).

\textbf{Variables derivadas} (calculadas en la consulta 10 mediante
\lstinline!$set!): \lstinline!frp_por_deteccion_mw! (FRP total dividido
entre detecciones medias) y \lstinline!proporcion_nominal_alta! (fracción de
detecciones de alta confianza).

\textbf{Extracción para análisis estadístico/ML:} la consulta 10 es
exactamente el mecanismo de extracción: un pipeline
\lstinline!$match+$set+$project! que produce un dataset plano —una fila por
celda-mes con columnas escalares— exportable directamente a
\texttt{pandas.DataFrame} para entrenar modelos o calcular estadísticas
descriptivas.

\section{Conclusiones}

\begin{enumerate}[itemsep=1pt]
  \item El modelo documental (objetos anidados + arreglo \texttt{evidencias}
  de cardinalidad fija) fue suficiente para responder las 10 consultas y
  los 2 pipelines analíticos exigidos sin necesidad de colecciones
  adicionales ni \lstinline!$lookup!.
  \item Los metadatos operativos del laboratorio se incorporaron de forma
  incremental (\lstinline!$set!, \lstinline!$push!) sin afectar ningún valor
  satelital, demostrando que el esquema soporta evolución posterior a la
  carga.
  \item Los hallazgos analíticos (estacionalidad septiembre-octubre, 2023
  como año de mayor intensidad) son válidos \textbf{dentro de esta muestra
  estratificada} y sirven de prueba de concepto para las agregaciones que
  el Proyecto Integrador ejecutaría sobre la población completa de
  celdas-mes.
\end{enumerate}

\textbf{Relación con el Proyecto Integrador:} este taller reutiliza la
unidad de análisis celda-mes de \texttt{AGENTS.md} y demuestra el componente
MongoDB de la arquitectura híbrida SQL--NoSQL de FireForest; el mismo
patrón de pipeline (consulta 10) alimentaría, en el proyecto completo, el
dataset analítico celda-mes junto con las tablas de hechos de
PostgreSQL/PostGIS (\texttt{fact\_incendio}, \texttt{fact\_clima}).

\textbf{Limitación del muestreo estratificado:} la muestra fija
deliberadamente 10 casos con fuego y 5 sin fuego por año (semilla 2026), por
lo que la proporción observada (70/35, 66,7\,\% con fuego) \textbf{no estima
la prevalencia real de incendios en el cantón Loja} y no debe generalizarse
a la población completa de celdas-mes. Ver
\texttt{03\_datos/README.md} para el detalle del método de muestreo.

\end{document}
